% ---------------------------------------------------------------------------
% Research-agenda addition (tenants). Inserted by the lead at the end of the
% conservation section, immediately before the canaries section.
% ---------------------------------------------------------------------------
\subsection{The ledger keyed by data subject}\label{sec:sealed-person-ledger}
\pdstatus{\Designed}{specified here; the reference system's A3 ledger keeps one balance}

\noindent The ledger above keeps one balance. A harbor (the book's name for
the governed environment in which agents work) that trains on many people's
logs needs one balance per data subject, because the guarantee each of them
cares about is their own. A \pdterm{subject-keyed ledger} keeps one meter per
data subject and charges each meter for every training round that could read
that subject's records. Such a harbor usually serves several customer organizations, each a
\pdterm{tenant} with its own agents, data and users. Two logs are
\pdterm{subject-level neighbors} when one is the other with every record of a
single data subject removed. This is a stronger relation than the one-record
neighbors of the $\varepsilon$-ledger, whose corollary protects a single
record, not a person with many records. Each record carries an owner label, here the set
of data subjects whose data it carries (Section~\ref{sec:sealed-tenants} gives the
full rule). The state is the chapter's, with one index added. Let $U$ be the
set of data subjects, and fix an order $\alpha>1$ of
\pdterm[Renyi differential privacy]{R\'enyi differential privacy} (RDP): a mechanism
is $(\alpha,\rho)$-RDP if the R\'enyi divergence of order $\alpha$ between its
output distributions on any two neighbors, in either direction, is at most
$\rho$ \pdcite{feldman2020renyi}. The R\'enyi divergence of order $\alpha$ of
$P$ from $Q$ is $\frac{1}{\alpha-1}\ln\mathbb{E}_{x\sim Q}\bigl[(P(x)/Q(x))^{\alpha}\bigr]$;
it is $0$ when $P=Q$ and grows as the two distributions separate. Each training round $t$ declares a
per-subject cost $\rho_t\ge0$, chosen from public information and previous
outputs.

\begin{pdclaim}{Theorem}{Subject-Keyed Conservation}
\setlength{\parskip}{4pt plus 1pt}
\label{thm:sealed-person-ledger}
Initially $\sigma(u)=0$ for every $u\in U$ and $\Lambda=()$, with budget
$B\ge0$. The only spending transition is $\mathrm{train}(t,\rho_t)$. Let
$O_t$ be the data subjects named in the owner labels of the round's candidate
records. The round computes the participating set
\[
P_t=\{u\in O_t:\sigma(u)+\rho_t\le B\},
\]
uses only records whose owner labels lie inside $P_t$, and \emph{atomically}
appends $(t,P_t,\rho_t,\text{artifact hash},\text{policy hash})$ to $\Lambda$
and sets $\sigma(u)\mathrel{+}=\rho_t$ for every $u\in P_t$.

Then every reachable state satisfies, for every data subject $u$,
\[
\sigma(u)=\sum_{e\in\Lambda,\;u\in P_e}\rho_e
\quad\text{and}\quad
\sigma(u)\le B,
\]
including under concurrent invocation, by the same single-writer
serialization as Theorem~\ref{thm:sealed-conservation}.

\textbf{Conditional privacy corollary.} Fix the subject-level neighboring
relation. Suppose every round is
$(\alpha,\rho_t)$-RDP for each data subject in $P_t$, conditional on the prior
transcript, with mechanisms and costs selected from public information and
previous outputs. Suppose also that the complete observer-visible interaction
is covered, including metadata, refusals and timing. The individual privacy
filter of Feldman and Zrnic \pdcite{feldman2020renyi} then gives every data subject
$(\alpha,B)$-RDP for the whole sequence of released artifacts, hence
$\bigl(B+\ln(1/\delta)/(\alpha-1),\,\delta\bigr)$-privacy for every
$\delta\in(0,1)$.
\end{pdclaim}
\pdscope{The accounting proof below establishes neither premise of the
corollary.}

\paragraph{Proof.}
Induction over the sequential interleaving, one coordinate per data subject. Base:
$\sigma\equiv0$ and $\Lambda=()$, so every sum is empty and every bound holds.
Step: a round appends one entry with set $P_t$ and cost $\rho_t$.
For $u\in P_t$ both sides of the equality grow by $\rho_t$, and the guard
$\sigma(u)+\rho_t\le B$ preserves the bound. For $u\notin P_t$ neither side
changes. Non-spending transitions touch neither component. Concurrency reduces
to the sequential case as in Theorem~\ref{thm:sealed-conservation}. With a
single data subject, the theorem is that theorem with $\rho$ in place of
$\varepsilon$. \qed

\begin{pdexample}{Four Meters, Three Rounds}
\setlength{\parskip}{4pt plus 1pt}
This is a constructed example. Take the budget $B=0.25$. Carol, Heron and Wren
have logs from the start. Round~1 costs $\rho_1=0.15$; all three can afford
it, so $P_1=\{\text{Carol},\text{Heron},\text{Wren}\}$ and each of their
meters reads $0.15$. Pike's first logs then arrive, and his meter starts
at~$0$.

Round~2 costs $0.10$, and everyone can afford it. Carol, Heron and Wren reach
$0.15+0.10=0.25$, exactly $B$, and Pike reaches $0.10$.

Round~3 costs $0.10$. Carol, Heron and Wren would reach $0.35>0.25$, so they are
out and $P_3=\{\text{Pike}\}$. A support thread labelled
$\{\text{Carol},\text{Pike}\}$ is now excluded, because Carol is out, even
though Pike is in; a record labelled $\{\text{Pike}\}$ is used. Pike ends
at $0.20$.

The journal holds three entries: round~1 for Carol, Heron and Wren at $0.15$;
round~2 for all four at $0.10$; round~3 for Pike alone at $0.10$. Each meter
equals the sum of the entries that name that subject: Carol
$0.15+0.10=0.25$, Pike $0.10+0.10=0.20$.

Now suppose Heron appears unresolved as two identities in two tenants. He then
has two meters, each capped at $0.25$, and can be charged $0.50$ in total.
\end{pdexample}

A round charges only the people whose records it could read. Removing a data subject
with no candidate record leaves the round's input unchanged, so the round costs
that subject nothing, and their meter does not move. This is what makes the
meters differ: in the example, Carol, Heron and Wren reach $B$ after two
rounds, while Pike, whose logs arrived later, stands at $0.20$. A ledger that charged
everyone with budget left would be sound but would move every meter in
lockstep.

The journal records $P_t$, and whether a data subject's records were used is itself
personal data. The journal therefore stays inside the confidential domain;
observers see its hash.

\paragraph{The Whitehouse warning, revisited.}
Section~\ref{sec:sealed-conservation} warns: ``Fully adaptive advanced
accounting would require a separately justified filter, such as Whitehouse et
al.'s \pdcite{wrrw23}, which matches advanced-composition rates and leading constants under
its stopping rule. That approximate-DP filter is not implemented by A3's
sum-only ledger.'' The warning stands for a ledger kept in $\varepsilon$.
Measured in R\'enyi terms, the sum-only rule is itself a valid filter under
full adaptivity: if adaptively chosen costs satisfy $\sum_t\rho_t\le B$ almost
surely, the composition is $(\alpha,B)$-RDP \pdcite{feldman2020renyi}.
Feldman and Zrnic report that the $(\varepsilon,\delta)$ filter this implies
is as tight as advanced composition obtained through concentrated
differential privacy. A3's ledger keeps its structure and needs only to
count in $\rho$ instead of $\varepsilon$.

\paragraph{Whose key.}
The guarantee holds for the key the ledger indexes. A data subject who appears
unresolved in $k$ tenants holds $k$ keys and can be charged $kB$. A data
subject named only inside other people's records holds no key at all, which is why
Table~\ref{tab:sealed-tenant-policy} keeps third parties out of training.

\begin{pdexample}{A Lifetime Budget of One Quarter}
\setlength{\parskip}{4pt plus 1pt}
\pdterm[zero-concentrated privacy]{Zero-concentrated privacy} with parameter $\rho$ means $(\alpha,\alpha\rho)$-RDP
for every $\alpha>1$, and for it the filter's rule is $\sum_t\rho_t\le\rho_{\max}$,
whatever the order \pdcite{feldman2020renyi}. Give each data subject
$\rho_{\max}=0.25$ and train in rounds of $\rho_t=0.025$: a subject joins
\[
0.25/0.025=10
\]
rounds and then drops out of training.

Convert with $(\alpha,\rho')$-RDP $\Rightarrow$
$\bigl(\rho'+\ln(1/\delta)/(\alpha-1),\delta\bigr)$, where here
$\rho'=\alpha\rho$. Write $x=\alpha-1$: then
$\varepsilon=\rho+x\rho+\ln(1/\delta)/x$, smallest at
$x=\sqrt{\ln(1/\delta)/\rho}$, which gives
\[
\varepsilon=\rho+2\sqrt{\rho\ln(1/\delta)}.
\]
At $\delta=10^{-9}$: $\ln10^9=20.723$, $0.25\times20.723=5.181$,
$\sqrt{5.181}=2.276$, so $\varepsilon=0.25+4.552=4.80$, at
$\alpha=1+\sqrt{82.89}\approx10.1$. At $\delta=10^{-10}$: $\ln10^{10}=23.026$,
$0.25\times23.026=5.756$, $\sqrt{5.756}=2.399$, so
$\varepsilon=0.25+4.799=5.05$.

Gboard's production language models, trained with private federated learning,
report $\rho$ between $0.2$ and $2$ at the device level, and with tighter
privacy-loss-distribution accounting quote $\varepsilon\approx4.49$ at
$\delta=10^{-10}$ for $\rho=0.25$ \pdcite{xu2023gboard}. The difference
$5.05-4.49=0.56$ comes from the simpler conversion; the guarantee itself is
the same.

\emph{Now you try:} rounds of $0.04$ against the same budget, at
$\delta=10^{-6}$ (Exercise~\ref{ex:sealed-person-budget}).
\end{pdexample}

\pdexercisepointerone{\ref{ex:sealed-person-budget}}{\pageref{ex:sealed-person-budget}}
